\documentclass[12pt, letterpaper]{article}

%% Build from the paper/ directory:
%%   cd paper && pdflatex main.tex
%% or use the root build system:
%%   make paper          (from repo root)
%%   ./build.sh paper    (POSIX)
%%   build.cmd paper     (Windows cmd)

% packages.tex
\usepackage{amsmath, amssymb, amsthm}
\usepackage{graphicx}
\usepackage{xcolor}
\usepackage[
    hidelinks,
    pdftitle={dcl-generator-zoo: A Catalogue of the 71-Dimensional Per-Site Automorphism Algebra of the A=1 Discrete Causal Lattice},
    pdfauthor={Jack Menendez},
    pdfsubject={Catalogue of the 71-dim per-site automorphism algebra; A=1 Discrete Causal Lattice; Paper~II Phase~3 companion},
    pdfkeywords={A=1 Discrete Causal Lattice, generator zoo, automorphism centralizer, SU(6), Standard Model gauge group, leptoquark, chirality, bracket structure, broken symmetry, sympy verification}
]{hyperref}
\usepackage{booktabs}
\usepackage{array}
\usepackage{geometry}
\usepackage{adjustbox}
\usepackage{longtable}
\usepackage{setspace}
\usepackage{caption}
\usepackage{subcaption}
\usepackage{float}
\onehalfspacing

\input{macros/commands}

%% Page geometry (also set here as a fallback in case the package path fails)
\geometry{top=1in, bottom=1in, left=1in, right=1in}

%% ---------------------------------------------------------------
%% TITLE PAGE
%% ---------------------------------------------------------------

\title{The dcl-generator-zoo:\\
       A Catalogue of the 71-Dimensional Per-Site\\
       Automorphism Algebra of the A=1 Discrete Causal Lattice%
       \thanks{Version 0.1-DRAFT.  Companion artifact to the A=1
       Discrete Causal Lattice series; built on the Phase~3
       centralizer enumeration of
       \emph{Geometry Forces Physics} (Paper~II,
       \href{https://doi.org/10.5281/zenodo.20240736}{DOI:
       10.5281/zenodo.20240736}), which itself builds on
       \emph{Geometry First} (Paper~I,
       \href{https://doi.org/10.5281/zenodo.20078529}{DOI:
       10.5281/zenodo.20078529}).  This catalogue is intended as
       the canonical reference for downstream papers in the series
       and for external work building on the framework's per-site
       algebraic structure.  DOI for this catalogue to be assigned
       at v1.0 deposit.}}
\author{Jack Menendez\\
        \small Eugene, OR\\
        \small ORCID: \href{https://orcid.org/0009-0003-1166-307X}{0009-0003-1166-307X}\\
        \small \href{mailto:jackdmenendez@gmail.com}{jackdmenendez@gmail.com}}
\date{\today}

\begin{document}

\maketitle

\begin{abstract}
    % abstract.tex

Phase~3 of Paper~II of the A=1 Discrete Causal Lattice series
established that the discrete-Hermitian centralizer of the bipartite
tick rule on the per-site $\mathbb{C}^{12}$ amplitude has dimension
71, structurally $\mathfrak{su}(6)_+ \oplus \mathfrak{su}(6)_-
\oplus \mathfrak{u}(1)$.  The 12-dimensional Standard Model gauge
subalgebra sits inside this centralizer as a non-normal subgroup,
and the 59-generator complement (``extras'') transforms under
$SU(3) \times SU(2)$ as a sum of $(\mathbf{8}, \mathbf{3})$
leptoquark-flavoured copies and chirality-twisted SM-flavoured
shadows.  This document catalogues all 71 generators: it gives each
one a stable name, tags its $(SU(3), SU(2))$ irrep content, records
its tensor-factor action on the chirality~$\otimes$~isospin~$\otimes$~colour
decomposition of $\mathbb{C}^{12}$, and reports its bracket class
against every other generator in the centralizer.  The catalogue is
emitted as machine-readable JSON for downstream consumption and as
a printable longtable for human reference.  The intent is that
follow-on papers in the series, and external researchers building
on the framework, can refer to specific generators by name rather
than re-deriving them.

\end{abstract}

%% ---------------------------------------------------------------
%% FRONT-MATTER NOTICES (after abstract per author preference)
%% ---------------------------------------------------------------

\medskip
\noindent\textbf{Keywords.}  A=1 Discrete Causal Lattice; per-site
automorphism algebra; centralizer; $\mathfrak{su}(6)$; Standard
Model gauge group; leptoquark; chirality; bracket structure; broken
symmetry; sympy verification.

\medskip
\noindent\textbf{License.}  Paper text and figures: Creative Commons
Attribution 4.0 (CC BY 4.0).  Source code in the companion
repository: MIT License.

\medskip
\noindent\textbf{Data availability.}  The catalogue, its
verification scripts, and the \LaTeX{} source for this document are
publicly available at the GitHub repository
\nolinkurl{github.com/JackDMenendez/dcl-generator-zoo}.  The
load-bearing centralizer enumeration scripts
\nolinkurl{automorphism_centralizer_extended.py} and
\nolinkurl{aut_centralizer_extras_commutators.py} were introduced
in Paper~II~\cite{menendez2026sm} and are reproduced here verbatim
as inherited evidence; the zoo's own enumeration script
\nolinkurl{generator_zoo.py} (planned for v1.0) extends them with
stable naming, irrep tagging, and JSON / \LaTeX{} rendering.

%% ---------------------------------------------------------------
%% RELEASE NOTES (only at release time)
%% ---------------------------------------------------------------
%% \bigskip
%% \noindent\textbf{Release notes for vX.Y (since vX.Y-1).}
%% \begin{itemize}
%% \item Item 1.
%% \end{itemize}

%% \emph{Status legend used throughout.}  \texttt{PASS} / \texttt{PART}
%% / \texttt{STUB} / \texttt{FAIL} -- see Table~\ref{tab:audit}.

\newpage
\tableofcontents
\newpage

%% ============================================================
%% MAIN
%% ============================================================
\section{Introduction}
% introduction.tex

\label{sec:introduction}

This document is the companion catalogue to Paper~II of the A=1
Discrete Causal Lattice series~\cite{menendez2026sm}, which
established Eq.~(137) of Paper~I~\cite{menendez2026geometry} as a
structural-equality conjecture between the per-site automorphism
algebra of the bipartite octahedral lattice and the Lie algebra
$\mathfrak{so}(3,1) \oplus \mathfrak{su}(3) \oplus \mathfrak{su}(2)
\oplus \mathfrak{u}(1)$.  Paper~II's Phase~3 enumeration produced
the sharper result that the discrete-Hermitian centralizer of the
bipartite tick rule on the per-site $\mathbb{C}^{12}$ amplitude has
dimension~71, with structural decomposition $\mathfrak{su}(6)_+
\oplus \mathfrak{su}(6)_- \oplus \mathfrak{u}(1)$.  The
12-dimensional Standard Model gauge subalgebra is contained in
this centralizer; the remaining 59 generators are the
``extras''~\cite{georgi1999lie}, structurally $(\mathbf{8},
\mathbf{3})$ leptoquark-flavoured plus chirality-$\sigma_x$-twisted
SM-flavoured.

That count -- 12 SM, 59 extras, 71 total -- is the input to this
catalogue.  Paper~II's prose treats the 59 extras at the level of
their representation-theoretic content and their bracket structure;
it does not give them individual names, irrep tags per generator,
or a machine-readable form that downstream papers can consume
without re-running the enumeration.  This zoo fills that gap.

\subsection{Scope}
\label{subsec:scope}

What this document is:

\begin{itemize}
\item A \emph{catalogue}, not a paper of new results.  All
      algebraic content is inherited from Paper~II Phase~3.
\item A persistent \emph{naming convention} for each of the 71
      centralizer generators.  Names are stable across catalogue
      revisions so that follow-on papers can cite generators by
      name.
\item A \emph{taxonomy} along several orthogonal axes (source
      $\mathfrak{su}(6)$ factor; $(SU(3), SU(2))$ irrep content;
      tensor-factor action; bracket class against every other
      generator).
\item A \emph{machine-readable emit} (JSON) plus a printable
      longtable rendering, so external work can consume the
      catalogue either programmatically or by reading.
\end{itemize}

What this document is not:

\begin{itemize}
\item A re-derivation of Paper~II Phase~3.  The dimensional and
      structural results are cited, not re-proved here.
\item A statement of physical interpretation.  The
      decoherence-channel hypothesis of Paper~II's
      \texttt{notes/extras\_as\_a1\_accounting.md} (DRAFT) is
      reproduced as an appendix note, marked clearly as
      conjectural; the catalogue itself records only the algebraic
      facts.
\item A successor to Paper~II.  This is companion infrastructure,
      not a Paper~III.
\end{itemize}

Table~\ref{tab:audit} records the catalogue's audit-row status.
Most rows inherit \texttt{PASS} from Paper~II Phase~3; the rows
specific to the zoo's value-add (stable naming, irrep tagging,
JSON emit, LaTeX renderer) are \texttt{STUB} in the v0.1 draft and
will flip to \texttt{PART}/\texttt{PASS} as Phase~2 of this
catalogue is executed.

%% Audit table embedded here so the reader meets it before any
%% claim that depends on a status entry.
% ============================================================
% audit_table.tex -- dcl-generator-zoo
%
% Catalogue audit: which catalogue features are PASS / PART / STUB.
%
% Rows split into two groups:
%   (1) "Inherited" rows whose algebraic content is verified by the
%       Paper~II Phase~3 scripts reproduced verbatim in
%       src/utilities/ (these are PASS in this zoo because they
%       were PASS in Paper~II).
%   (2) "Zoo-specific" rows whose evidence is new in this catalogue
%       (stable naming, irrep tagging, JSON emit, LaTeX renderer,
%       bracket-class lookup, conjectural decoherence-channel
%       labels).  These start STUB in v0.1 and flip as Phase~2
%       lands.
%
% Status semantics (paste into the abstract or front-matter as well):
%   PASS  -- derivation complete or script confirms the claim to
%            stated precision.
%   PART  -- mechanism demonstrated, full quantitative match pending;
%            specific gaps documented in the evidence column.
%   STUB  -- analytical result with no numerical confirmation yet,
%            or planned script not yet written.
%   FAIL  -- tested and disconfirmed.
% ============================================================

\label{tab:audit}

{\scriptsize
\setlength{\tabcolsep}{4pt}
\renewcommand{\arraystretch}{0.95}
\begin{longtable}{@{}p{3.4cm}p{4.0cm}p{4.6cm}p{1.0cm}@{}}
\caption{\small Catalogue audit for the dcl-generator-zoo.  Rows in
the first block (``inherited'') are algebraic results established by
Paper~II Phase~3 and reproduced here verbatim via the scripts under
\texttt{src/utilities/}.  Rows in the second block (``zoo-specific'')
are catalogue features new in this artifact.  Status semantics in the
front-matter (\texttt{PASS} / \texttt{PART} / \texttt{STUB} /
\texttt{FAIL}).}
\label{tab:audit_caption} \\
\hline
\textbf{Claim} & \textbf{Mechanism} & \textbf{Evidence} & \textbf{Status} \\
\hline
\endfirsthead

\multicolumn{4}{l}{\textit{(Table~\ref{tab:audit}, continued from previous page)}} \\
\hline
\textbf{Claim} & \textbf{Mechanism} & \textbf{Evidence} & \textbf{Status} \\
\hline
\endhead

\hline
\multicolumn{4}{r}{\textit{(continued on next page)}} \\
\endfoot

\hline
\endlastfoot

% ── INHERITED FROM PAPER II PHASE 3 ─────────────────────────────────
\multicolumn{4}{l}{\textit{Inherited from Paper~II~\cite{menendez2026sm} Phase~3 (reproduced verbatim).}} \\
\hline

Centralizer dimension is 71
    & Enumerate the 143-element Pauli\,$\otimes$\,Pauli\,$\otimes$\,Gell-Mann
      tensor basis of $\mathfrak{su}(12)$; count basis elements
      commuting with $\sigma_x \otimes I_2 \otimes I_3$
    & \texttt{src/utilities/automorphism\_centralizer\_extended.py}
      reports 71, matching $\dim(\mathfrak{u}(6) \oplus
      \mathfrak{u}(6)) - 1$
    & \texttt{PASS} \\

SM gauge subalgebra has dim 12 inside the centralizer
    & $1$ ($J_1$) $+\,3$ ($\mathfrak{su}(2)_W$) $+\,8$
      ($\mathfrak{su}(3)_c$); verify each is in the centralizer
    & Same script, Step~2; 12 of 14 conjecture Hermitian generators
      are in the centralizer ($J_2, J_3$ are continuum-emergent)
    & \texttt{PASS} \\

Extras count is 59, with breakdown 24 + 3 + 8 + 24
    & Iso--colour mixings (24), chirality-$\sigma_x$ + isospin (3),
      chirality-$\sigma_x$ + colour (8), chirality-$\sigma_x$ +
      iso--colour (24)
    & Same script, Step~3; counts match predictions
    & \texttt{PASS} \\

71-dim centralizer is
$\mathfrak{su}(6)_+ \oplus \mathfrak{su}(6)_- \oplus \mathfrak{u}(1)$
    & $\sigma_x$-eigenbasis decomposition $\mathbb{C}^{12} =
      \mathbb{C}^6_+ \oplus \mathbb{C}^6_-$; centralizer is
      block-diagonal $\mathfrak{u}(6)_\pm$
    & \texttt{notes/aut\_centralizer\_enumeration.md}; structural
      reading is derivable directly from eigenspace dimensions
    & \texttt{PASS} \\

$(\mathbf{8}, \mathbf{3})$ + chirality-shadow SM
decomposition of the 59 extras under $SU(3) \times SU(2)$
    & $\mathbf{35} = (\mathbf{8}, \mathbf{1}) \oplus (\mathbf{1},
      \mathbf{3}) \oplus (\mathbf{8}, \mathbf{3})$ branching of
      $\mathfrak{su}(6)$ adjoint; applied to each chirality block
    & \texttt{notes/aut\_centralizer\_enumeration.md};
      48 leptoquark-flavoured + 11 chirality-twisted SM-flavoured =
      59
    & \texttt{PASS} \\

$[\text{extras}, \text{SM}]$ never lands in SM (extras is an
SM-invariant module)
    & Compute all 708 brackets symbolically; classify
      zero/in\_SM/in\_extras/mixed by rank in the SM basis
    & \texttt{src/utilities/aut\_centralizer\_extras\_commutators.py}
      reports 256 zero, 0 in\_SM, 452 in\_extras
    & \texttt{PASS} \\

$[\text{extras}, \text{extras}]$ has SM components (extras is not a
Lie ideal)
    & Compute all 1711 $\binom{59}{2}$ brackets; classify by
      span membership
    & Same script, Step~4b; 339 zero, 178 in\_SM, 1146 in\_extras,
      48 mixed
    & \texttt{PASS} \\

% ── ZOO-SPECIFIC ────────────────────────────────────────────────────
\hline
\multicolumn{4}{l}{\textit{Zoo-specific catalogue features (new in this artifact).}} \\
\hline

Stable name for each of the 71 generators
    & Persistent identifier (\texttt{G\_01} .. \texttt{G\_71})
      ordered by SM block first then extras categories; preserved
      across catalogue revisions
    & \texttt{src/utilities/generator\_zoo.py} (\texttt{build\_catalogue()});
      naming-order rationale in docstring; awaits user-side first run
    & \texttt{PART} \\

Tensor-factor 3-bit action tag (chirality, isospin, colour) per
generator
    & For each generator, record which of the three
      $\mathbb{C}^2 \otimes \mathbb{C}^2 \otimes \mathbb{C}^3$
      tensor factors it acts on non-trivially
    & \texttt{src/utilities/generator\_zoo.py}
      (\texttt{derive\_factor\_tag()}); derived from tensor signature
      already in \texttt{automorphism\_centralizer\_extended.py}
    & \texttt{PART} \\

$(SU(3), SU(2))$ irrep tag per generator
    & Irrep tag from the
      $\mathbf{35} = (\mathbf{8}, \mathbf{1}) + (\mathbf{1},
      \mathbf{3}) + (\mathbf{8}, \mathbf{3}) + (\mathbf{1},
      \mathbf{1})$ branching, projected per generator
    & \texttt{src/utilities/generator\_zoo.py}
      (\texttt{derive\_irrep()}); informally used in
      \texttt{notes/aut\_centralizer\_extras\_commutators.md}
    & \texttt{PART} \\

Bracket-class lookup table for every pair $(X, Y)$ of centralizer
generators
    & Extend the $59 \times 12$ + $\binom{59}{2}$ bracket
      classification of Paper~II to the full $\binom{71}{2} = 2485$
      table and emit a lookup
    & \texttt{src/utilities/generator\_zoo.py}
      (\texttt{classify\_all\_brackets()}); imports
      \texttt{rank\_of\_matrices} / \texttt{is\_in\_span} from the
      inherited Paper~II commutator script
    & \texttt{PART} \\

JSON catalogue emit for machine-readable downstream consumption
    & Generator zoo writes
      \texttt{data/generator\_catalogue.json} with one record per
      generator (name, tensor signature, factor tag, irrep tag,
      SM classification) plus the full bracket-class table
    & \texttt{src/utilities/generator\_zoo.py}
      (\texttt{emit\_json()}); schema version 1; JSON file is
      written on first run
    & \texttt{PART} \\

\LaTeX{} longtable renderer for the catalogue appendix
    & Generator zoo also emits a \LaTeX{} fragment that the paper's
      appendix \texttt{\textbackslash input}s; \texttt{IfFileExists}
      placeholder ensures the build succeeds even before the
      script's first run
    & \texttt{src/utilities/generator\_zoo.py}
      (\texttt{emit\_latex\_longtable()}); appendix section
      \texttt{paper/sections/generator\_zoo.tex} hand-written;
      the table proper \texttt{paper/sections/generator\_zoo\_table.tex}
      is auto-generated
    & \texttt{PART} \\

Conjectural decoherence-channel labels
    & Per-generator label flagging the
      decoherence-as-projection conjecture of
      \texttt{notes/extras\_as\_a1\_accounting.md}
    & Conjecture itself is \texttt{DRAFT} (see note); labels are
      out of scope for v0.1.  The per-generator records do carry
      a \texttt{sm\_class} field (\texttt{SM} / \texttt{extras})
      and a \texttt{sm\_subgroup} field which a downstream
      decoherence-channel consumer can read; the *physical-role*
      labels are deferred per \texttt{feedback\_let\_formalism\_mature}
    & \texttt{STUB} \\

\hline
\end{longtable}
}



%% Future sections (Phase~2 of this catalogue):
%%   \section{Naming convention}
%%   \section{Tensor-factor action tags}
%%   \section{Representation-theoretic content}
%%   \section{Bracket-class lookup}
%%   \section{JSON emit and downstream consumption}
%%   \section{Generator zoo} (the catalogue appendix proper)

\section{Conclusion}
% conclusion.tex

\placeholder{Conclusion to be written in Phase~2 once the generator
zoo's stable naming, irrep tagging, JSON emit, and \LaTeX{} renderer
are in place.  At v0.1-DRAFT this section is intentionally a
placeholder; the substantive concluding content depends on what the
catalogue actually contains.}


\section*{Acknowledgements}
% acknowledgements.tex

\placeholder{Acknowledge collaborators, reviewers, and tools.  The
inherited centralizer enumeration scripts were introduced in
Paper~II of the A=1 Discrete Causal Lattice series and rely on the
sympy Lie-algebra and matrix infrastructure.  Acknowledgement of
arXiv endorsers (if any) and Zenodo community curators (the
\texttt{a1-discrete-causal-lattice}, \texttt{fbt-framework}, and
\texttt{osqgr} communities at minimum) belongs here once the
catalogue is deposited.}


%% ============================================================
%% APPENDICES
%% ============================================================
\appendix

\section{Generator Zoo}
% generator_zoo.tex -- the catalogue appendix.
%
% The longtable proper is auto-generated by
% src/utilities/generator_zoo.py from data/generator_catalogue.json
% and \input{}'d at the end of this section.  This file provides
% the hand-written reader's map, per-class structural notes, and
% pointers to the JSON schema.

\label{sec:generator_zoo}

The 71 generators of the per-site automorphism centralizer are
catalogued in Table~\ref{tab:generator_zoo}.  This appendix
describes the row schema, the seven SM-classification blocks the
table is partitioned into, and the JSON emit that machine-readable
downstream consumers should use rather than scraping the
\LaTeX{} table.

\subsection{Row schema}
\label{subsec:zoo_row_schema}

Each of the 71 rows of Table~\ref{tab:generator_zoo} records:

\begin{description}
\item[Name.]  A stable identifier of the form \texttt{G\_NN} with
  \texttt{NN} $\in \{01, \ldots, 71\}$.  Names are preserved across
  catalogue revisions and are the canonical reference for downstream
  papers.  The ordering convention is documented in the docstring
  of \texttt{src/utilities/generator\_zoo.py}; in brief: SM block
  first (\texttt{G\_01} = $J_1$, \texttt{G\_02--G\_04} = $SU(2)_W$,
  \texttt{G\_05--G\_12} = $SU(3)_c$), then extras in the order
  iso-colour mixings, $\sigma_x$-iso, $\sigma_x$-colour,
  $\sigma_x$-iso-col.
\item[Chir / Iso / Col.]  The tensor-signature factor of the
  generator on $\mathbb{C}^{12} = \mathbb{C}^2_{\text{chir}}
  \otimes \mathbb{C}^2_{\text{iso}} \otimes \mathbb{C}^3_{\text{col}}$.
  Possible values: $I$ (identity), $\sigma_x / \sigma_y / \sigma_z$
  (Paulis, for the chirality and isospin factors), or
  $\lambda_1, \ldots, \lambda_8$ (Gell-Manns, for the colour
  factor).  Only $I$ and $\sigma_x$ appear in the chirality column
  because the centralizer condition forces $[X, \sigma_x \otimes
  I_2 \otimes I_3] = 0$.
\item[Tag.]  A 3-bit factor-action tag $(c, i, k) \in \{0, 1\}^3$,
  one bit per factor, with bit value $1$ iff the corresponding
  tensor factor is non-identity.  Quick at-a-glance reading of
  which tensor factors a generator touches.
\item[$\mathbf{SU(3)}$.]  The colour irrep, either
  $\mathbf{8}$ (adjoint of $SU(3)_c$) or $\mathbf{1}$ (trivial).
\item[$\mathbf{SU(2)}$.]  The weak-isospin irrep, either
  $\mathbf{3}$ (adjoint of $SU(2)_W$) or $\mathbf{1}$ (trivial).
\item[Block.]  The chirality eigenspace ($I$ or $\sigma_x$) that
  the generator lives in.  In the structural decomposition
  $\mathfrak{Aut}_{\text{cent}} = \mathfrak{su}(6)_+ \oplus
  \mathfrak{su}(6)_- \oplus \mathfrak{u}(1)$, the $I$-block
  generators sit in $\mathfrak{su}(6)_+$ and the $\sigma_x$-block
  generators in $\mathfrak{su}(6)_-$ (with $J_1$ as the central
  $\mathfrak{u}(1)$).
\item[Class.]  The SM-classification subgroup the generator
  belongs to.  Seven possible values: \texttt{J\_1},
  \texttt{SU(2)\_W}, \texttt{SU(3)\_c} (the 12 SM generators);
  \texttt{iso\_col\_mixing} (24, leptoquark-flavoured),
  \texttt{sigmaX\_iso} (3, chirality-shadow $SU(2)_W$),
  \texttt{sigmaX\_col} (8, chirality-shadow $SU(3)_c$),
  \texttt{sigmaX\_iso\_col} (24, chirality-shadow leptoquark).
\end{description}

\subsection{The seven blocks at a glance}
\label{subsec:zoo_blocks}

\begin{description}
\item[\texttt{J\_1} (1 generator).]  The central $\mathfrak{u}(1)$
  of the centralizer.  $J_1 = \sigma_x \otimes I_2 \otimes I_3$
  commutes with everything in the 71-dim algebra and is the
  hypercharge-like central charge of the SM block.

\item[\texttt{SU(2)\_W} (3 generators).]  The standard weak-isospin
  generators $I_2 \otimes \sigma_a \otimes I_3$, $a \in \{x, y, z\}$.
  These are SM gauge.

\item[\texttt{SU(3)\_c} (8 generators).]  The standard colour
  Gell-Manns $I_2 \otimes I_2 \otimes \lambda_a$, $a \in \{1,
  \ldots, 8\}$.  These are SM gauge.

\item[\texttt{iso\_col\_mixing} (24 generators).]  Leptoquark-
  flavoured operators $I_2 \otimes \sigma_a \otimes \lambda_b$.
  Transform as $(\mathbf{8}, \mathbf{3})$ under the SM adjoint
  action -- colour-octet, weak-triplet, exactly the GUT
  leptoquark irrep.  In the chirality-$I$ block.

\item[\texttt{sigmaX\_iso} (3 generators).]  Chirality-shadow
  $SU(2)_W$ operators $\sigma_x \otimes \sigma_a \otimes I_3$.
  Same $(\mathbf{1}, \mathbf{3})$ SM-irrep content as
  $\mathfrak{su}(2)_W$, but living in the $\sigma_x$ chirality
  block, so they are not SM gauge.

\item[\texttt{sigmaX\_col} (8 generators).]  Chirality-shadow
  $SU(3)_c$ operators $\sigma_x \otimes I_2 \otimes \lambda_a$.
  Same $(\mathbf{8}, \mathbf{1})$ content as $\mathfrak{su}(3)_c$,
  in the $\sigma_x$ block.

\item[\texttt{sigmaX\_iso\_col} (24 generators).]  Chirality-shadow
  leptoquark operators $\sigma_x \otimes \sigma_a \otimes \lambda_b$.
  Same $(\mathbf{8}, \mathbf{3})$ content as the
  \texttt{iso\_col\_mixing} block, in the $\sigma_x$ block.
\end{description}

Sums: SM block $= 1 + 3 + 8 = 12$; extras block $= 24 + 3 + 8 + 24
= 59$; total $= 71$.  Of the 59 extras, $48 = 24 + 24$ are
leptoquark-flavoured (in $(\mathbf{8}, \mathbf{3})$) and the
remaining $11 = 3 + 8$ are chirality-shadow SM-flavoured (in
$(\mathbf{1}, \mathbf{3}) \oplus (\mathbf{8}, \mathbf{1})$).

\subsection{Machine-readable emit}
\label{subsec:zoo_json}

The catalogue is also emitted as
\texttt{data/generator\_catalogue.json}, with one record per
generator (same fields as Table~\ref{tab:generator_zoo}) plus a
full bracket-class table: for each of the
$\binom{71}{2} = 2485$ unordered pairs $(G_i, G_j)$ with $i < j$,
the class of the bracket $[G_i, G_j]$ (one of \texttt{zero} /
\texttt{in\_SM} / \texttt{in\_extras} / \texttt{mixed}) is
recorded.  Downstream tools that consume the catalogue
programmatically should read the JSON and refer to generators
by their \texttt{G\_NN} name; this LaTeX longtable is the human-
facing companion.

The JSON schema is documented in the docstring of
\texttt{src/utilities/generator\_zoo.py}.  Regenerate the JSON
and this longtable in lockstep via:

\begin{verbatim}
python -m src.utilities.generator_zoo
\end{verbatim}

\subsection{Bracket-class summary}
\label{subsec:zoo_brackets}

The full bracket-class table is in the JSON emit; here we
record the aggregate counts (also reproduced by
\texttt{src/utilities/aut\_centralizer\_extras\_commutators.py}
for the subsets it covers):

\begin{itemize}
\item $[\text{SM}, \text{SM}]$ (66 pairs): closure of the 12-dim
  SM gauge subalgebra inside the centralizer.  All brackets land
  in SM (the SM block is a Lie subalgebra).
\item $[\text{extras}, \text{SM}]$ (708 pairs): 256 zero, 0
  in\_SM, 452 in\_extras, 0 mixed.  The 59 extras form an
  invariant module under the adjoint action of SM.
\item $[\text{extras}, \text{extras}]$ (1711 pairs): 339 zero,
  178 in\_SM, 1146 in\_extras, 48 mixed.  The extras subspace is
  \emph{not} a Lie ideal; brackets within it produce SM
  components, consistent with the standard algebraic shape of a
  broken-symmetry pattern (see
  \texttt{notes/aut\_centralizer\_extras\_commutators.md}).
\end{itemize}

Sum: $66 + 708 + 1711 = 2485 = \binom{71}{2}$.

\subsection{The catalogue}
\label{subsec:zoo_catalogue}

%% The longtable below is auto-generated by
%% src/utilities/generator_zoo.py.  Do not edit by hand; regenerate
%% via 'python -m src.utilities.generator_zoo'.  Wrapped in
%% \IfFileExists{} so a fresh clone whose user has not yet run
%% the catalogue script still builds; the placeholder branch then
%% prints a visible reminder.
\IfFileExists{sections/generator_zoo_table.tex}{%
    \input{sections/generator_zoo_table}%
}{%
    \placeholder{Auto-generated catalogue table not present.
    Regenerate via \texttt{python -m src.utilities.generator\_zoo}
    from the repo root; the script writes
    \texttt{paper/sections/generator\_zoo\_table.tex} and
    \texttt{data/generator\_catalogue.json} together.}%
}


\section{Code and Data}
% code_and_data.tex
\label{sec:code_and_data}

\subsection{Repository layout}

The catalogue's source tree is laid out as follows:

\begin{itemize}
\item \texttt{src/utilities/automorphism\_centralizer\_extended.py}
  -- Paper~II Phase~3 centralizer enumeration, reproduced verbatim.
  Builds the 143-element Pauli\,$\otimes$\,Pauli\,$\otimes$\,Gell-Mann
  tensor basis of $\mathfrak{su}(12)$, filters to the 71-dim
  centralizer of the bipartite tick rule, classifies into 7
  tensor-signature categories.
\item \texttt{src/utilities/aut\_centralizer\_extras\_commutators.py}
  -- Paper~II Phase~3 bracket-classification script, reproduced
  verbatim.  Computes all 708 $[\text{extras}, \text{SM}]$ brackets
  and 1711 $[\text{extras}, \text{extras}]$ brackets, classifies
  each as zero / in\_SM / in\_extras / mixed.
\item \texttt{src/utilities/generator\_zoo.py} (planned for
  v1.0) -- the catalogue script proper.  Extends the two
  inherited scripts above with stable per-generator names, an
  $(SU(3), SU(2))$ irrep tag, a 3-bit tensor-factor action tag,
  and a bracket-class lookup; emits a JSON catalogue and a
  \LaTeX{} longtable fragment.
\item \texttt{data/generator\_catalogue.json} (planned for v1.0)
  -- machine-readable emit of the catalogue, one record per
  generator.
\item \texttt{notes/} -- four Paper~II notes reproduced with
  Provenance prefixes, providing the textual companion to the
  inherited scripts.
\item \texttt{paper/sections/generator\_zoo.tex} (planned for v1.0)
  -- the catalogue appendix proper (longtable + reader's map).
\end{itemize}

\subsection{Reproducing the v0.1 audit table}

Every \texttt{PASS} row in Table~\ref{tab:audit} is reproduced by:

\begin{verbatim}
python -u src/utilities/automorphism_centralizer_extended.py
python -u src/utilities/aut_centralizer_extras_commutators.py
\end{verbatim}

The two scripts run independently; each prints a structured report
of the centralizer dimension, the SM-aligned vs extras split, and
the bracket-classification counts.  Total runtime is on the order
of minutes on a modern laptop (the second script computes 2419
symbolic $12 \times 12$ commutators and classifies each by rank).


\section{Reproducibility Quickstart}
% reproducibility.tex
\label{sec:reproducibility}

\begin{itemize}
\item \textbf{Environment.}  Python~3 with \texttt{sympy},
      \texttt{numpy}, \texttt{matplotlib}, \texttt{pytest}.
      Exact dependency list in \texttt{virtual-env-requirements.txt}.
      The \texttt{setup.sh} / \texttt{setup.cmd} scripts create a
      \texttt{.venv} and install everything.

\item \textbf{Build the paper.}  From a fresh clone, with the venv
      activated:
      \begin{verbatim}
./build.sh paper        # POSIX / MSYS2 UCRT64 on Windows
build.cmd paper         # Windows cmd / PowerShell
      \end{verbatim}
      runs \texttt{pdflatex} three times and \texttt{bibtex} once,
      producing \texttt{paper/main.pdf}.

\item \textbf{Verify the headline result.}  The two inherited
      Paper~II scripts together reproduce every \texttt{PASS} row
      in Table~\ref{tab:audit}:
      \begin{verbatim}
python -u src/utilities/automorphism_centralizer_extended.py
python -u src/utilities/aut_centralizer_extras_commutators.py
      \end{verbatim}
      Each prints a structured report ending in either ``\texttt{PASS}''
      or a count-mismatch error.  Total runtime: a few minutes.

\item \textbf{Run the master audit roll-up.}  From the repo root:
      \begin{verbatim}
python audit_universe.py
      \end{verbatim}
      parses \texttt{paper/sections/audit\_table.tex} and reports
      the declared status of every row.  Rows whose evidence is
      analytical (script output cited as
      \texttt{src/utilities/*.py}) are reported as-is.

\item \textbf{Where to look next.}  The four notes under
      \texttt{notes/} (each prefixed with a Provenance block)
      provide the textual companion to the algebraic results.
      Start with \texttt{notes/aut\_centralizer\_enumeration.md}
      for the structural identification of the 71-dim centralizer
      as $\mathfrak{su}(6)_+ \oplus \mathfrak{su}(6)_- \oplus
      \mathfrak{u}(1)$; then
      \texttt{notes/aut\_centralizer\_extras\_commutators.md}
      for the bracket-classification reading; then (DRAFT)
      \texttt{notes/extras\_as\_a1\_accounting.md} for the
      conjectural decoherence-as-projection physical
      interpretation.
\end{itemize}


%% ============================================================
\bibliographystyle{unsrt}
\bibliography{paper-bib/references}

\end{document}
